\documentclass[
    language=english,
    doctype=dissertation-proposal,
    institution=none,
]{../../omnilatex}

\RequirePackage{config/document-settings}

\begin{document}
\maketitle

\begin{abstract}
This dissertation proposal outlines a research program investigating
\textit{Probabilistic Abstractions for Scalable Bayesian Inference in
High-Dimensional Models}. The proposed research addresses a fundamental
tension in modern machine learning: the need for principled uncertainty
quantification (provided by Bayesian methods) versus the computational
intractability of exact Bayesian inference in large-scale models.

The dissertation develops three interconnected research thrusts: (1) a
theoretical framework for constructing provably accurate probabilistic
abstractions of complex models, (2) scalable inference algorithms that
leverage these abstractions to achieve sublinear computational complexity,
and (3) empirical validation on real-world problems in climate science
and computational biology where uncertainty quantification is critical.

The expected contributions include new theory bridging information geometry
and approximate inference, open-source software for scalable Bayesian
computation, and applications demonstrating the practical utility of the
proposed methods.
\end{abstract}

\section{Introduction and Motivation}
\label{sec:introduction}

\subsection{Problem Statement}

Bayesian inference provides a principled framework for incorporating prior
knowledge, quantifying uncertainty, and performing model selection. However,
exact Bayesian inference is intractable for all but the simplest models.
Modern machine learning systems --- with millions to billions of parameters
--- present an extreme challenge for Bayesian methods.

Current approximate inference techniques (variational inference, MCMC, Laplace
approximations) face a fundamental scalability bottleneck:

\begin{equation}
    \text{Cost} \propto \mathcal{O}(nd^2 + d^3)
    \label{eq:bottleneck}
\end{equation}

where $n$ is the number of data points and $d$ is the number of parameters.
For models with $d > 10^6$, these costs become prohibitive.

\subsection{Research Vision}

This dissertation proposes that \textit{probabilistic abstractions} ---
compressed representations of model components that preserve essential
statistical properties --- can break the scalability barrier while
maintaining rigorous uncertainty guarantees.

\subsection{Central Hypothesis}

\begin{quote}
    \textit{Probabilistic abstractions, defined as information-geometric
    projections of model components onto tractable parametric families,
    can reduce the computational complexity of Bayesian inference from
    $\mathcal{O}(nd^2 + d^3)$ to $\mathcal{O}(nkd + k^3)$, where $k \ll d$
    is the abstraction dimension, while maintaining approximation error
    bounds that are controlled by the information geometry of the
    projection.}
\end{quote}

\section{Research Questions}
\label{sec:research-questions}

This dissertation addresses the following research questions:

\begin{enumerate}
    \item \textbf{RQ1 (Theory):} How can we construct probabilistic
          abstractions that minimize information loss while maximizing
          computational savings? What are the fundamental limits of
          abstraction accuracy?

    \item \textbf{RQ2 (Algorithms):} Can we develop inference algorithms
          that operate directly on abstractions, achieving sublinear
          complexity in the number of parameters while maintaining
          calibrated uncertainty estimates?

    \item \textbf{RQ3 (Applications):} Do probabilistic abstractions
          enable Bayesian inference on previously intractable problems
          in climate science and computational biology, and do the
          resulting uncertainty estimates improve downstream decision-making?
\end{enumerate}

\section{Literature Review}
\label{sec:literature}

\subsection{Approximate Bayesian Inference}

The literature on approximate Bayesian inference spans several decades and
multiple communities.

\begin{table}[htbp]
    \centering
    \caption{Summary of approximate inference methods and their limitations.}
    \label{tab:inference-methods}
    \begin{tabular}{@{}lccp{5cm}@{}}
        \toprule
        Method & Complexity & Scalability & Key Limitation \\
        \midrule
        MCMC & $\mathcal{O}(nd^2)$ per sample & Poor & Mixing in high dimensions \\
        Variational Inference & $\mathcal{O}(ndk)$ & Good & Mean-field independence assumption \\
        Laplace Approximation & $\mathcal{O}(nd^2 + d^3)$ & Poor & Mode-seeking, underestimates variance \\
        Expectation Propagation & $\mathcal{O}(ndk)$ & Moderate & Non-convex optimization, no convergence guarantees \\
        MCMC + subsampling & $\mathcal{O}(n' d^2)$ & Moderate & Gradient noise from minibatch \\
        \bottomrule
    \end{tabular}
\end{table}

\subsection{Model Compression and Abstraction}

Model compression techniques (pruning, quantization, knowledge distillation)
have been developed primarily for deterministic models. Their Bayesian
counterparts --- which must preserve uncertainty information --- are less
well understood. Recent work on Bayesian compression \cite{blundell2015weight}
and probabilistic pruning provides a foundation, but does not address the
abstraction problem directly.

\subsection{Information Geometry and Inference}

Information geometry \cite{amari2016information} provides a natural
framework for analyzing approximate inference. The Fisher information
metric defines a Riemannian manifold on the space of probability
distributions, and variational inference can be interpreted as a
geodesic projection onto a tractable family.

\begin{equation}
    D_{\text{KL}}(p \| q) = \int p(x) \log \frac{p(x)}{q(x)} \, dx
    \label{eq:kl}
\end{equation}

The KL divergence is not symmetric, and its geometry determines which
approximation family is most suitable for a given model component.

\section{Proposed Research}
\label{sec:proposed}

\subsection{Thrust 1: Theoretical Framework for Probabilistic Abstractions}

\subsubsection{Objective}
Develop a rigorous mathematical framework for constructing, analyzing, and
composing probabilistic abstractions of model components.

\subsubsection{Approach}
\begin{enumerate}
    \item Define abstractions as information-geometric projections: given a
          complex distribution $p$, find $q^* \in \mathcal{Q}$ (a tractable
          family) that minimizes $D_{\text{KL}}(p \| q)$ or $D_{\text{KL}}(q \| p)$.
    \item Derive error bounds for downstream inference tasks that depend on
          the abstraction quality.
    \item Develop composition theorems: how do abstraction errors propagate
          when multiple components are abstracted?
\end{enumerate}

\subsubsection{Expected Outcome}
A set of theorems characterizing the fundamental trade-off between
abstraction dimension $k$ and inference accuracy, analogous to the
bias-variance trade-off but in the information-geometric setting.

\subsection{Thrust 2: Scalable Inference Algorithms}

\subsubsection{Objective}
Develop inference algorithms that operate on probabilistic abstractions to
achieve sublinear computational complexity.

\subsubsection{Approach}
\begin{enumerate}
    \item Design a hierarchical variational inference algorithm where model
          layers are represented at different abstraction levels.
    \item Develop adaptive abstraction: automatically determine the
          appropriate abstraction dimension for each model component based
          on its contribution to posterior uncertainty.
    \item Implement GPU-accelerated inference kernels for the proposed
          algorithms.
\end{enumerate}

\subsubsection{Deliverable}
An open-source software package (\texttt{ProbAbst.jl}) implementing the
proposed algorithms in Julia, with a Python API.

\subsection{Thrust 3: Applications}

\subsubsection{Objective}
Demonstrate the practical utility of probabilistic abstractions on two
real-world domains where uncertainty quantification is critical.

\subsubsection{Application 1: Climate Model Emulation}

Climate models (GCMs) are computationally expensive, making Bayesian model
emulation challenging. We will develop probabilistic abstractions of GCM
components to enable uncertainty-aware climate projections.

\subsubsection{Application 2: Single-Cell RNA Sequencing}

Single-cell RNA-seq datasets contain millions of cells and tens of thousands
of genes. Current Bayesian methods for cell-type annotation and trajectory
inference are limited to small datasets. We will scale these methods using
probabilistic abstractions.

\section{Methodology}
\label{sec:methodology}

\subsection{Theoretical Methods}

\begin{itemize}
    \item Information geometry and Riemannian optimization
    \item PAC-Bayesian generalization bounds
    \item Transport map theory for probability distributions
    \item Functional analysis of kernel operators
\end{itemize}

\subsection{Computational Methods}

\begin{itemize}
    \item Variational inference with structured variational families
    \item Automatic differentiation (Julia Zygote, Python JAX)
    \item GPU kernel optimization (CUDA, ROCm)
    \item Distributed computing for large-scale experiments
\end{itemize}

\subsection{Empirical Methods}

\begin{table}[htbp]
    \centering
    \caption{Planned experimental evaluations.}
    \label{tab:experiments}
    \begin{tabular}{@{}lccc@{}}
        \toprule
        Experiment & Dataset & Baseline & Metrics \\
        \midrule
        Abstraction quality & Synthetic & VI, MCMC & KL divergence, coverage \\
        Scalability & Synthetic & VI, MCMC & Wall-clock time, memory \\
        Climate emulation & CMIP6 & GPE, Deep GP & RMSE, calibration \\
        scRNA-seq & 10X Genomics & SCVI, PAGA & ARI, NMI \\
        \bottomrule
    \end{tabular}
\end{table}

\section{Timeline}
\label{sec:timeline}

\begin{table}[htbp]
    \centering
    \caption{Research timeline and milestones.}
    \label{tab:timeline}
    \begin{tabular}{@{}lcd{12}l@{}}
        \toprule
        Period & Task & {Milestone} \\
        \midrule
        Year 1, Q1--Q2 & Literature review and theory development & Thrust 1 draft \\
        Year 1, Q3--Q4 & Error bounds and composition theorems & Conference paper \\
        Year 2, Q1--Q2 & Inference algorithm design & Algorithm prototype \\
        Year 2, Q3--Q4 & Software implementation & ProbAbst.jl v0.1 \\
        Year 3, Q1--Q2 & Climate emulation application & Domain paper \\
        Year 3, Q3--Q4 & scRNA-seq application & Domain paper \\
        Year 4, Q1--Q2 & Dissertation writing & Thesis draft \\
        Year 4, Q3--Q4 & Defense preparation & Dissertation defense \\
        \bottomrule
    \end{tabular}
\end{table}

\section{Expected Contributions}
\label{sec:contributions}

\begin{enumerate}
    \item \textbf{Theoretical:} A formal framework for probabilistic
          abstractions with provable accuracy guarantees.
    \item \textbf{Algorithmic:} Scalable inference algorithms achieving
          $\mathcal{O}(nkd + k^3)$ complexity.
    \item \textbf{Software:} Open-source package \texttt{ProbAbst.jl}
          for Bayesian inference with probabilistic abstractions.
    \item \textbf{Applications:} Demonstrated utility in climate science
          and computational biology.
    \item \textbf{Publications:} At least 3 conference papers (NeurIPS,
          ICML, AISTATS) and 2 journal papers (JMLR, Biostatistics).
\end{enumerate}

\section{Broader Impacts}
\label{sec:broader-impacts}

\subsection{Scientific Impact}

Scalable Bayesian inference with uncertainty quantification has the
potential to improve scientific reasoning across disciplines, from climate
policy (where uncertainty bounds on projections are critical for
decision-making) to genomics (where single-cell heterogeneity must be
quantified, not merely summarized).

\subsection{Societal Impact}

By enabling uncertainty-aware machine learning at scale, this research
could improve the reliability and trustworthiness of AI systems used in
high-stakes domains such as healthcare, environmental policy, and
autonomous systems.

\subsection{Broader Implications}

The information-geometric framework developed here may have applications
beyond machine learning, including signal processing, computational
neuroscience, and statistical physics, where compressed representations
of complex systems are essential.

\section{Conclusion}
\label{sec:conclusion}

This dissertation proposes a research program that addresses a fundamental
challenge in modern machine learning: principled Bayesian inference at
scale. Through the development of probabilistic abstractions, the proposed
research aims to provide both theoretical foundations and practical tools
for uncertainty-aware machine learning in high-dimensional settings.

\bibliographystyle{plain}
\bibliography{references}

\end{document}
